\documentclass[10pt,a4paper]{amsart}
\usepackage{amsmath,amssymb,amsthm,graphicx}
% Retained preamble for the Stacks Project chapter {CH} (shared project preamble).
% The chapter file pinned-original-source.tex only carries \input{{preamble}} on its first line,
% so no preamble text lives inside the chapter itself.  The Stacks Project uses one single shared
% preamble for every chapter: the snapshot stages it three times (sources/stacks-chow/preamble.tex,
% sources/stacks-etale-cohomology/preamble.tex and focused-seeds/AGP/sources/stacks-snapshot-20261011/
% preamble.tex) and those copies are identical in size.  The block below is copied mechanically,
% line for line, from sources/stacks-chow/preamble.tex lines 14-18 and 154-217: the Xy-pic setup and
% the native theorem declarations together with the mathematical macros.  Nothing was added, removed
% or redefined; the wrapper class and the amsmath/amssymb/amsthm/graphicx base packages are supplied
% by the assembler.
% For commutative diagrams we use Xy-pic
% Chapter: stacks-cohomology.tex -- Cohomology of Sheaves.
\usepackage[all]{xy}

% We use 2cell for 2-commutative diagrams.
\xyoption{2cell}

% Theorem environments.
%
\theoremstyle{plain}
\newtheorem{theorem}[subsection]{Theorem}
\newtheorem{proposition}[subsection]{Proposition}
\newtheorem{lemma}[subsection]{Lemma}

\theoremstyle{definition}
\newtheorem{definition}[subsection]{Definition}
\newtheorem{example}[subsection]{Example}
\newtheorem{exercise}[subsection]{Exercise}
\newtheorem{situation}[subsection]{Situation}

\theoremstyle{remark}
\newtheorem{remark}[subsection]{Remark}
\newtheorem{remarks}[subsection]{Remarks}

\numberwithin{equation}{subsection}

% Macros
%
\def\lim{\mathop{\mathrm{lim}}\nolimits}
\def\colim{\mathop{\mathrm{colim}}\nolimits}
\def\Spec{\mathop{\mathrm{Spec}}}
\def\Hom{\mathop{\mathrm{Hom}}\nolimits}
\def\Ext{\mathop{\mathrm{Ext}}\nolimits}
\def\SheafHom{\mathop{\mathcal{H}\!\mathit{om}}\nolimits}
\def\SheafExt{\mathop{\mathcal{E}\!\mathit{xt}}\nolimits}
\def\Sch{\mathit{Sch}}
\def\Mor{\mathop{\mathrm{Mor}}\nolimits}
\def\Ob{\mathop{\mathrm{Ob}}\nolimits}
\def\Sh{\mathop{\mathit{Sh}}\nolimits}
\def\NL{\mathop{N\!L}\nolimits}
\def\CH{\mathop{\mathrm{CH}}\nolimits}
\def\proetale{{pro\text{-}\acute{e}tale}}
\def\etale{{\acute{e}tale}}
\def\QCoh{\mathit{QCoh}}
\def\Ker{\mathop{\mathrm{Ker}}}
\def\Im{\mathop{\mathrm{Im}}}
\def\Coker{\mathop{\mathrm{Coker}}}
\def\Coim{\mathop{\mathrm{Coim}}}

% Boxtimes
%
\DeclareMathSymbol{\boxtimes}{\mathbin}{AMSa}{"02}

%
% Macros for moduli stacks/spaces
%
\def\QCohstack{\mathcal{QC}\!\mathit{oh}}
\def\Cohstack{\mathcal{C}\!\mathit{oh}}
\def\Spacesstack{\mathcal{S}\!\mathit{paces}}
\def\Quotfunctor{\mathrm{Quot}}
\def\Hilbfunctor{\mathrm{Hilb}}
\def\Curvesstack{\mathcal{C}\!\mathit{urves}}
\def\Polarizedstack{\mathcal{P}\!\mathit{olarized}}
\def\Complexesstack{\mathcal{C}\!\mathit{omplexes}}
% \Pic is the operator that assigns to X its picard group, usage \Pic(X)
% \Picardstack_{X/B} denotes the Picard stack of X over B
% \Picardfunctor_{X/B} denotes the Picard functor of X over B
\def\Pic{\mathop{\mathrm{Pic}}\nolimits}
\def\Picardstack{\mathcal{P}\!\mathit{ic}}
\def\Picardfunctor{\mathrm{Pic}}
\def\Deformationcategory{\mathcal{D}\!\mathit{ef}}
\usepackage{cleveref}
\crefname{originallemma}{original lemma}{original lemmas}
\Crefname{originallemma}{Original lemma}{Original lemmas}
\makeatletter
\newlabel{modules-lemma-abelian}{{01AG}{source}{}{sourcepointer-1}{}}
\newlabel{modules-lemma-abelian@cref}{{[originallemma][0][]01AG}{[1][0][]source}{}{}{}}
\makeatother
\crefname{originaldefinition}{original definition}{original definitions}
\Crefname{originaldefinition}{Original definition}{Original definitions}
\makeatletter
\newlabel{more-algebra-definition-K-flat}{{06XZ}{source}{}{sourcepointer-2}{}}
\newlabel{more-algebra-definition-K-flat@cref}{{[originaldefinition][0][]06XZ}{[1][0][]source}{}{}{}}
\makeatother
\title{K-flatness of a complex of modules is checked on stalks}
\author{The Stacks Project Authors}
\date{stacks-project:cohomology.tex:local-snapshot-20261011; GNU FDL 1.2 or later}
\begin{document}
\maketitle
\noindent\textit{Editorial note.} All delivered mathematical text is the author's own, copied from the pinned Stacks Project chapter file cohomology.tex; no statement or proof step was rewritten, reordered or removed.\par\medskip
\noindent\textit{Editorial note.} The proof is complete as written and ends with the stalk computation; no step is referred to the reader.\par\medskip
\noindent\textit{Editorial note.} The counted spans carry 2 navigation stubs and no bibliography citation.\par\medskip
\noindent\textit{Editorial note.} The preamble is the chapter's retained minimal preamble (sources/stacks-cohomology/retained-preamble.tex lines 1-81), a verbatim excerpt of the native preamble.\par\medskip
\noindent\textit{Editorial note.} The label lemma-check-K-flat-stalks also occurs in sites-cohomology.tex. The ruling recorded for this lane before assembly is to keep this card: the sites statement is about a complex whose terms are flat modules over a ringed topos and its conclusion includes flatness of the terms, so the two are related by a genuinely stronger hypothesis and a genuinely stronger conclusion rather than being one statement restated.\par\medskip
\noindent\textit{Allowed standard prerequisites (editorial).}\begin{itemize}
\item Exactness of complexes of sheaves is checked on stalks. Modules, Lemma modules-lemma-abelian, invoked at proof lines 5885-5886. Supplies the first direction.
\item Definition of K-flatness for complexes over a ring. More on Algebra, Definition more-algebra-definition-K-flat, named in the statement at line 5878. Supplies the property tested on stalks.
\item Skyscraper pushforward along a point and its interaction with stalks and tensor products. Standard sheaf theory, used at proof lines 5889-5893. Supplies the acyclic complex of sheaves used to transfer the K-flat property from the space to the local ring.
\end{itemize}\medskip
\ifdefined\hypertarget\hypertarget{sourcepointer-1}{}\fi
\noindent\textit{Source reference.} Original lemma 01AG: sources/stacks-chow/tags.txt, source lines 1689--1689. Pointer outside the counted span, invoked once for context; the load-bearing argument of the counted statement is the author's proof reproduced in full in the proof span.\par\medskip
\ifdefined\hypertarget\hypertarget{sourcepointer-2}{}\fi
\noindent\textit{Source reference.} Original definition 06XZ: sources/stacks-chow/tags.txt, source lines 8868--8868. Pointer outside the counted span, invoked once for context; the load-bearing argument of the counted statement is the author's proof reproduced in full in the proof span.\par\medskip
\begin{lemma}
\label{lemma-check-K-flat-stalks}
Let $(X, \mathcal{O}_X)$ be a ringed space. Let $\mathcal{K}^\bullet$
be a complex of $\mathcal{O}_X$-modules. Then $\mathcal{K}^\bullet$
is K-flat if and only if for all $x \in X$ the complex
$\mathcal{K}_x^\bullet$ of $\mathcal{O}_{X, x}$-modules is K-flat
(More on Algebra, Definition \ref{more-algebra-definition-K-flat}).
\end{lemma}

\begin{proof}
If $\mathcal{K}_x^\bullet$ is K-flat for all $x \in X$ then we see
that $\mathcal{K}^\bullet$ is K-flat because $\otimes$ and
direct sums commute with taking stalks and because we can check exactness
at stalks, see
Modules, Lemma \ref{modules-lemma-abelian}.
Conversely, assume $\mathcal{K}^\bullet$ is K-flat. Pick $x \in X$
and let
$M^\bullet$ be an acyclic complex of $\mathcal{O}_{X, x}$-modules.
Then $i_{x, *}M^\bullet$ is an acyclic complex of $\mathcal{O}_X$-modules.
Thus $\text{Tot}(i_{x, *}M^\bullet \otimes_{\mathcal{O}_X} \mathcal{K}^\bullet)$
is acyclic. Taking stalks at $x$ shows that
$\text{Tot}(M^\bullet \otimes_{\mathcal{O}_{X, x}} \mathcal{K}_x^\bullet)$
is acyclic.
\end{proof}

\end{document}
